\section*{Резюме}

Настоящата дипломна работа разглежда проектирането, реализацията и оценяването на библиотека за обектно-релационно съпоставяне по време на компилация за C++23~\cite{cpp23standard}. Работата е организирана около две допълващи се оси. Първата е моделирането на данните и заявките по време на компилация, при което логическата структура на достъпа до данни не се описва чрез низове, интерпретирани по време на изпълнение, а чрез статично типизирани \code{constexpr} конструкции. Втората е базираната на корутини асинхронна архитектура, чрез която същият типово ограничен модел се разширява към неблокиращо или псевдо-неблокиращо изпълнение според възможностите на конкретния драйвер на свързващия компонент.

Разработената библиотека използва статичен модел на данните, изразен в C++ чрез конструкции като \code{property<T, Name>}, \code{relationship<...>} и \code{table\_name\_trait<T>}, както и типизирано междинно представяне на заявкитете, реализирано чрез \code{select\_query}, \code{insert\_query}, \code{update\_query} и \code{delete\_query}. Този логически слой е независим от конкретния тип база данни. Материализацията към релационни и нерелационни системи се осъществява чрез специализации на \code{connector\_trait<DB>}, които декларират поддържаните възможности и превеждат една и съща логическа заявка към Structured Заявка Language (SQL), Binary JSON (BSON)-подобни филтри, Redis команди, Cypher или Cassandra Заявка Language (CQL) според избрания свързващ компонент.

Съществено разширение на архитектурата е асинхронният слой, изграден върху корутини, \code{Task<T>}, \code{thread\_pool}, \code{async\_db<DB>} и управление на връзки, съобразено с корутини. Показано е, че единна асинхронна абстракция може да бъде реализирана по два начина: чрез изнасяне на блокиращи операции към набор от нишки и чрез платформено-специфични неблокиращи интерфейси за драйвери, които ги предоставят. Така работата не само описва логическа ОРС абстракция, но и изследва изпълнителния модел, чрез който тази абстракция остава практически използваема при конкурентни натоварвания.

В дипломната работа се анализират както релационните сценарии за SQLite, PostgreSQL и MySQL, така и нерелационните сценарии за MongoDB, Redis, Neo4j и Cassandra. Показано е как един и същ C++ модел диктува логическата схема, но се материализира различно като таблица, документ, ключ-стойност структура, графов възел или широк колонен ред. Отделно се разглеждат клиентските библиотеки на съответните платформи и техните семантики на изпълнение, за да се определи къде библиотеката може честно да предложи платформено-специфично асинхронно поведение и къде правилният избор остава синхронен драйвер с контролирано изнесено изпълнение.

Верификацията на разработката се основава на модулни и интеграционни тестове от хранилището, а емпиричната оценка се допълва от сценарии за измерване на производителността за сравнение между синхронния и асинхронния режим. Те покриват изграждането на заявките, ограничаването на операции въз основа на възможностите, подготвените заявки, миграциите, нишковата безопасност, инфраструктурата за корутини и поведението на конкретните конектори при работа с реални или имитиращи драйвери. Това позволява формулиране не само на архитектурен модел, но и на конкретни критерии за „напълно имплементиран и функциониращ“ конектор и за практически оправдан асинхронен слой.

Основният принос на дипломната работа е в четири направления. Първо, предложен е практически използваем модел за библиотека за обектно-релационно съпоставяне по време на компилация, който съчетава статична типова безопасност и архитектура, независима от свързващия компонент. Второ, показано е как C++ описанието на данните може да служи като първичен източник за логическа схема при различни типове бази данни. Трето, изградена е базирана на корутини асинхронна архитектура, която разширява модела на конектора без компромис с ограниченията по време на компилация. Четвърто, формализиран е договор на конектора, чрез който библиотеката може да се надгражда с нови реализации на свързващи компоненти и асинхронни възможности без промяна в потребителския програмен интерфейс за заявки.

\vspace{1cm}

\textbf{Ключови думи:} C++23, обектно-релационно съпоставяне по време на компилация, статична типова безопасност, корутини, асинхронно изпълнение, \code{connector\_trait}, релационни и нерелационни бази данни, емпирична оценка
